% 数理物理1 数理物理1-12 問題1（092-bunsan。当時は 数理物理1-10 問題1・shou10-l1-bunsan） 解いた日 2026-10-03
%% kotae
\toi{1}{◎}{}
（図：質点を輪状に並べた図）
\begin{align*}
  \omega_\mu&=\sqrt{\frac{2K}{m}\cdot2\sin^{2}\Bigl(\frac{\pi\mu}{N}\Bigr)}\\
  &=2\sqrt{\frac{K}{m}}\sin\Bigl(\frac{\pi\mu}{N}\Bigr)\\
  &=2\sqrt{\frac{K}{m}}\abs*{\sin\Bigl(\frac{ak_\mu}{2}\Bigr)}
\end{align*}
（脇に $k_\mu=2\pi\mu/(Na)$，$ak_\mu/2=\pi\mu/N$）
\toi{2}{◎}{}
\[ =2\sqrt{\frac{K}{m}}\cdot\frac{ak_\mu}{2}=\underbrace{a\sqrt{\frac{K}{m}}}_{v_s}k_\mu \]
\toi{3}{×}{未着手}
\mikai
\toi{4}{×}{未着手}
\mikai
%% kizuki
\item (3)(4) は手を付けていない。(3) は先に $\sqrt{K/m}$ の値を出し，$v_s=a\sqrt{K/m}$ と $\omega_{\max}=2\sqrt{K/m}$（$\abs{\sin}=1$ となる $k_\mu a=\pi$，$\mu=N/2$）を有効数字2桁と単位つきで答える。
\item (4) は $x=k_\mu a$ とおくと比が $\omega_\mu/(v_sk_\mu)=2\sin(x/2)/x$ と $x$ だけの関数になる。$\sin(x/2)\simeq x/2-x^{3}/48$ で $1-x^{2}/24$ とし，波長は $\lambda_\mu=2\pi a/x$ で求める。
